\documentclass[10pt,a4paper]{article}
\usepackage{bayan}
% Figures for the Bayan final report. Every number plotted is measured; the
% source of each series is named in the caption that uses it.

\tikzset{
  flow/.style={-{Stealth[length=4pt]}, draw=black!80, line width=0.6pt},
  built/.style={draw=sicblue, line width=0.8pt, fill=white, rounded corners=2pt,
                minimum height=9mm, align=center, inner xsep=4pt},
  chosen/.style={built, fill=figpale, dash pattern=on 3pt off 1.5pt},
  pending/.style={draw=figgrey, line width=0.7pt, dashed, fill=white,
                  rounded corners=2pt, minimum height=9mm, align=center,
                  inner xsep=4pt, text=black!70},
  gate/.style={draw=figwarm, line width=0.9pt, fill=figpale, rounded corners=2pt,
               minimum height=9mm, align=center, inner xsep=4pt},
}

% ------------------------------------------------------------- schedule --
% Gantt of the real timeline. 13 Sep is day 0; 4.3mm per day.
\newcommand{\figschedule}{%
\begin{tikzpicture}[font=\scriptsize, x=4.3mm, y=-4.6mm]
  \def\lab{-0.6}
  % day grid and date labels
  \foreach \d/\t in {0/13,2/15,4/17,6/19,8/21,10/23,12/25,14/27,16/29,18/1,20/3,22/5} {
    \draw[black!12] (\d,0.3) -- (\d,16.9);
    \node[anchor=south, text=black!60] at (\d,0.3) {\t};
  }
  \node[anchor=south east, text=black!60] at (\lab,0.3) {September \enspace$|$\enspace day};
  \node[anchor=south, text=black!60] at (19.5,-0.35) {October};
  % bars: row, start, end, style, label
  \foreach \r/\s/\e/\c/\t in {
    1/0/5/sicblue/{Scoping, literature audit, brief},
    2/4/6/sicblue/{Repository and review workflow \#1--\#3},
    3/6/8/sicblue/{Product-form decision \#9},
    4/6/9/sicblue/{Validators, readability check \#10 \#17},
    5/6/9/sicblue/{Evaluation scripts, locked tests \#7},
    6/6/9/sicblue/{Diacritization benchmark \#6},
    7/6/10/sicblue/{SAMER conversion, model 1 \#5},
    8/7/9/sicblue/{Android plugin shell \#14},
    9/8/10/sicblue/{Copy baselines, converter \#19},
    10/8/10/sicblue/{Model 1 diagnosis; Baseet filtered},
    11/5/9/sicblue/{Action Plan}} {
    \fill[\c] (\s,\r-0.32) rectangle (\e,\r+0.32);
    \node[anchor=east] at (\lab,\r) {\t};
  }
  % in progress: solid to today, hatched onward
  \foreach \r/\s/\e/\t in {12/6/16/{Synthetic corpus \#8 \#12}, 13/10/16/{Model 2: mixed corpus, AraBART}} {
    \fill[figwarm] (\s,\r-0.3) rectangle (10,\r+0.3);
    \draw[figwarm, line width=0.7pt, fill=figwarm!12] (10,\r-0.3) rectangle (\e,\r+0.3);
    \node[anchor=east] at (\lab,\r) {\t};
  }
  \foreach \r/\s/\e/\t in {14/11/19/{On-device inference, int8}, 15/15/20/{Human evaluation}, 16/18/23/{Final report, slides, pitch}} {
    \draw[figgrey, dashed, line width=0.6pt, fill=black!5] (\s,\r-0.3) rectangle (\e,\r+0.3);
    \node[anchor=east] at (\lab,\r) {\t};
  }
  % milestones
  \foreach \d/\t/\a in {9/{AP}/north east, 10/{M1}/north west, 17/{M2}/north,
                         22/{Submit}/north east, 23/{Pitch}/north west} {
    \node[fill=black, diamond, inner sep=1.3pt] at (\d,17.4) {};
    \node[anchor=\a, font=\tiny] at (\d,17.7) {\t};
  }
  \node[anchor=east] at (\lab,17.4) {Milestones};
  % today
  \draw[figwarm, line width=1pt] (10,0.3) -- (10,16.9);
  \node[anchor=south, fill=white, inner sep=1pt, text=figwarm, font=\tiny\bfseries] at (10,16.9) {today};
\end{tikzpicture}}

% ---------------------------------------------------------- architecture --
\newcommand{\figarchitecture}{%
\begin{tikzpicture}[font=\footnotesize]
  \node[built, minimum width=30mm] (host) {Any app\\{\scriptsize message, web page, PDF}};
  \node[built, right=9mm of host, minimum width=32mm] (intent)
       {\texttt{PROCESS\_TEXT}\\{\scriptsize menu entry \ar{تبسيط}}};
  \node[built, right=9mm of intent, minimum width=30mm] (sheet)
       {Bottom sheet\\{\scriptsize simplified + original}};
  \draw[flow] (host) -- (intent);
  \draw[flow] (intent) -- (sheet);

  \node[pending, below=24mm of host.west, anchor=west, minimum width=24mm]
       (strip) {strip diacritics};
  \node[chosen, right=4mm of strip, minimum width=29mm] (simp)
       {simplifier\\{\scriptsize model 1 trained}};
  \node[chosen, right=4mm of simp, minimum width=25mm] (diac)
       {diacritizer\\{\scriptsize Libtashkeel}};
  \node[pending, right=4mm of diac, minimum width=23mm] (tts)
       {read-aloud\\{\scriptsize word timings}};
  \draw[flow] (strip) -- (simp);
  \draw[flow] (simp) -- (diac);
  \draw[flow] (diac) -- (tts);
  \begin{scope}[on background layer]
    \node[draw=figgrey, line width=0.6pt, rounded corners=3pt, inner xsep=6pt,
          inner ysep=6pt, fit=(strip)(simp)(diac)(tts)] (dev) {};
  \end{scope}
  \node[anchor=north west, font=\scriptsize\itshape, text=black!65]
       at ([yshift=-1pt]dev.south west) (devlab) {on the device, ONNX Runtime Mobile, no network};
  \draw[flow] (intent.south) -- ++(0,-6mm) -| (strip.north);
  \draw[flow] (tts.north) -- ++(0,6mm) -| (sheet.south);

  % legend
  \node[built, minimum height=3mm, minimum width=6mm, below=3mm of devlab.south west,
        anchor=north west] (k1) {};
  \node[right=1mm of k1, font=\scriptsize] (l1) {built and tested};
  \node[chosen, minimum height=3mm, minimum width=6mm, right=5mm of l1] (k2) {};
  \node[right=1mm of k2, font=\scriptsize] (l2) {selected or trained, not yet in the app};
  \node[pending, minimum height=3mm, minimum width=6mm, right=5mm of l2] (k3) {};
  \node[right=1mm of k3, font=\scriptsize] {not started};
\end{tikzpicture}}

% --------------------------------------------------------------- pipeline --
\newcommand{\figpipeline}{%
\begin{tikzpicture}[font=\footnotesize, node distance=8mm]
  \node[built, minimum width=34mm] (src) {source corpora\\{\scriptsize SAMER, BAREC, Baseet, DAASI}};
  \node[gate, below=of src, minimum width=34mm] (leak) {leakage check\\{\scriptsize exact + 5-gram Jaccard}};
  \node[built, below=of leak, minimum width=34mm] (train) {fine-tuning\\{\scriptsize model 1 on SAMER}};
  \node[built, below=of train, minimum width=34mm] (eval) {scoring\\{\scriptsize beside the copy baseline}};
  \draw[flow] (src) -- (leak);
  \draw[flow] (leak) -- node[right, font=\scriptsize, text=figwarm] {clean rows only} (train);
  \draw[flow] (train) -- (eval);

  \node[built, left=24mm of leak, minimum width=30mm, dash pattern=on 4pt off 2pt]
       (test) {locked test sets\\{\scriptsize SHA-256 manifest}};
  \draw[flow] (test) -- (leak);
  \draw[flow] (test.south) |- node[pos=0.75, below, font=\scriptsize] {final models only} (eval.west);

  \node[pending, right=24mm of src, minimum width=30mm] (gen) {LLM generation\\{\scriptsize synthetic pairs}};
  \node[gate, right=24mm of leak, minimum width=30mm] (judge) {validators\\{\scriptsize + human review}};
  \draw[flow, figgrey, dashed] (gen) -- (judge);
  \draw[flow, figgrey, dashed] (judge) |- (train.east);
\end{tikzpicture}}

% -------------------------------------------------------------------- EDA --
\newcommand{\figlengths}{%
\begin{tikzpicture}
\begin{axis}[
  width=\linewidth, height=55mm, font=\footnotesize,
  xlabel={source sentence length (words)}, ylabel={share of sentences (\%)},
  symbolic x coords={1--5,6--10,11--15,16--20,21--25,26--30,31--35,36--40,$>$40},
  xtick=data, ymin=0, ymax=45, grid=major, grid style={black!8},
  axis x line*=bottom, axis y line*=left,
  legend style={at={(0.98,0.97)}, anchor=north east, draw=black!25, font=\scriptsize},
  every axis plot/.append style={line width=1pt, mark size=1.4pt}]
  \addplot[sicblue, mark=*] coordinates {(1--5,37.17)(6--10,39.41)(11--15,15.58)(16--20,5.31)(21--25,1.49)(26--30,0.66)(31--35,0.15)(36--40,0.11)($>$40,0.10)};
  \addlegendentry{SAMER train (14,343)}
  \addplot[figgrey, mark=square*] coordinates {(1--5,29.05)(6--10,22.82)(11--15,17.28)(16--20,12.45)(21--25,7.80)(26--30,4.65)(31--35,2.54)(36--40,1.39)($>$40,2.01)};
  \addlegendentry{BAREC train (54,845)}
  \addplot[figwarm, mark=triangle*] coordinates {(1--5,11.25)(6--10,23.92)(11--15,18.08)(16--20,13.45)(21--25,10.36)(26--30,7.78)(31--35,5.40)(36--40,3.76)($>$40,6.00)};
  \addlegendentry{Baseet, filtered (31,104)}
  \addplot[black, dashed, mark=o] coordinates {(1--5,19.15)(6--10,25.38)(11--15,27.77)(16--20,18.09)(21--25,6.36)(26--30,2.12)(31--35,0.66)(36--40,0.46)($>$40,0.00)};
  \addlegendentry{DAASI (1,509)}
\end{axis}
\end{tikzpicture}}

% ---------------------------------------------------------- training run --
\newcommand{\figtraining}{%
\begin{tikzpicture}
\begin{axis}[
  width=\linewidth, height=56mm, font=\footnotesize,
  xlabel={training step}, ylabel={copy rate (\%)},
  xmin=100, xmax=3350, ymin=58, ymax=95,
  axis y line*=left, axis x line*=bottom,
  legend style={at={(0.02,0.03)}, anchor=south west, draw=black!25, font=\scriptsize},
  grid=major, grid style={black!8}]
  \addplot[figwarm, line width=1pt, mark=*, mark size=1.1pt] coordinates {
    (200,91.4)(400,90.2)(600,88.2)(800,85.6)(1000,77.8)(1200,75.2)(1400,68.8)
    (1600,71.8)(1800,70.2)(2000,69.0)(2200,68.6)(2400,67.6)(2600,66.8)
    (2800,63.6)(3000,65.2)(3200,64.8)(3250,64.8)};
  \addlegendentry{copy rate (left axis)}
  \addlegendimage{sicblue, line width=1pt, mark=square*, mark size=1.1pt}
  \addlegendentry{SARI (right axis)}
\end{axis}
\begin{axis}[
  width=\linewidth, height=56mm, font=\footnotesize,
  xmin=100, xmax=3350, ymin=71, ymax=77.5,
  axis y line*=right, axis x line=none, ylabel={SARI, 500-row dev subset}]
  \addplot[sicblue, line width=1pt, mark=square*, mark size=1.1pt] coordinates {
    (200,72.35)(400,73.69)(600,73.99)(800,74.13)(1000,73.75)(1200,74.17)
    (1400,74.34)(1600,74.87)(1800,75.27)(2000,74.90)(2200,75.29)(2400,75.14)
    (2600,75.61)(2800,75.42)(3000,75.74)(3200,75.56)(3250,75.56)};
  \draw[black!50, line width=0.6pt, dotted] (axis cs:2200,71) -- (axis cs:2200,77.5);
  \node[font=\scriptsize, text=black!60, anchor=north east, align=right]
    at (axis cs:2170,77.3) {validation loss\\bottoms here};
  \draw[black!50, line width=0.6pt, dotted] (axis cs:3000,71) -- (axis cs:3000,77.5);
  \node[font=\scriptsize, text=black!60, anchor=south east, align=right]
    at (axis cs:2970,71.2) {checkpoint\\selected};
\end{axis}
\end{tikzpicture}}

% The decoder and behaviour charts are generated with the numbers into results.tex.

% --------------------------------------------------------------------- UI --
\definecolor{appteal}{HTML}{0D5D56}
\definecolor{appsurf}{HTML}{FCFDF7}
\newcommand{\phone}[2]{% {name}{position}
  \node[draw=black!70, line width=1pt, rounded corners=5pt, minimum width=46mm,
        minimum height=86mm, fill=white] (#1) at #2 {};}
\newcommand{\figui}{%
\begin{tikzpicture}[font=\scriptsize]
  % (a) host app with the selection menu
  \phone{p1}{(0,0)}
  \foreach \y in {30,24,18,12,6,0,-6,-12,-18,-24,-30}
    \fill[black!10] ([xshift=-19mm, yshift=\y mm]p1.center) rectangle ++(38mm,2mm);
  \fill[sicblue!25] ([xshift=-12mm, yshift=6mm]p1.center) rectangle ++(28mm,2mm);
  \node[draw=black!30, fill=white, rounded corners=2pt, drop shadow={opacity=0.15},
        inner sep=3pt, font=\scriptsize] at ([yshift=15mm]p1.center)
       {\ar{نسخ}\quad\ar{تحديد الكل}\quad\textcolor{appteal}{\bfseries\ar{تبسيط}}};
  \node[below=2mm of p1, font=\scriptsize\bfseries] {(a) select text in any app};
  % (b) the bottom sheet
  \phone{p2}{(58mm,0)}
  \foreach \y in {30,24,18}
    \fill[black!10] ([xshift=-19mm, yshift=\y mm]p2.center) rectangle ++(38mm,2mm);
  \fill[black, opacity=0.35, rounded corners=5pt] (p2.north west) rectangle (p2.south east);
  \fill[appsurf, rounded corners=6pt] ([yshift=-18mm]p2.north west) rectangle (p2.south east);
  \fill[black!30, rounded corners=1pt] ([xshift=-3mm, yshift=-20mm]p2.north) rectangle ++(6mm,0.8mm);
  \node[anchor=north east, text=appteal, font=\scriptsize\bfseries] at ([xshift=-3mm, yshift=-22mm]p2.north east) {\ar{تبسيط بيان}};
  \node[draw=appteal, line width=0.8pt, rounded corners=3pt, minimum width=40mm,
        minimum height=19mm, anchor=north] (simp) at ([yshift=-28mm]p2.north) {};
  \node[anchor=north east, text=appteal, font=\tiny\bfseries] at ([xshift=-2mm, yshift=-1.5mm]simp.north east) {\ar{النص المبسط للقراءة الميسرة}};
  \foreach \y in {7,11,15}
    \fill[black!55] ([xshift=-2mm, yshift=-\y mm]simp.north east) rectangle ++(-30mm,1.4mm);
  \node[draw=black!25, fill=black!4, rounded corners=3pt, minimum width=40mm,
        minimum height=12mm, anchor=north] (orig) at ([yshift=-2mm]simp.south) {};
  \node[anchor=north east, text=black!55, font=\tiny] at ([xshift=-2mm, yshift=-1.2mm]orig.north east) {\ar{النص الأصلي}};
  \foreach \y in {6,9}
    \fill[black!25] ([xshift=-2mm, yshift=-\y mm]orig.north east) rectangle ++(-34mm,1mm);
  \node[anchor=north, font=\tiny, text=appteal] at ([yshift=-3mm]orig.south)
       {\ar{استمع}\qquad\ar{نسخ}\qquad\ar{إعدادات}};
  \node[below=2mm of p2, font=\scriptsize\bfseries] {(b) result sheet over the app};
  % typography callout
  \node[anchor=west, align=left, font=\scriptsize, text width=34mm] (t) at ([xshift=6mm, yshift=4mm]p2.east)
       {20\,sp text\\line height 1.6 + 10\,sp\\letter spacing 0.04\\right-to-left, \emph{not} justified};
  \draw[black!50, -{Stealth[length=3pt]}] (t.west) -- ([xshift=1mm]simp.east);
\end{tikzpicture}}

% Generated by make_results.py -- do not edit by hand.
\newcommand{\copytestall}{77.50}
\newcommand{\copytestchanged}{56.06}
\newcommand{\testchangedn}{1,678}
\newcommand{\testunchangedn}{1,599}
\newcommand{\greedyall}{75.88}
\newcommand{\greedychanged}{61.83}
\newcommand{\greedyunchanged}{90.62}
\newcommand{\beamall}{76.29}
\newcommand{\beamchanged}{61.95}
\newcommand{\beamunchanged}{91.34}
\newcommand{\greedygain}{+5.77}
\newcommand{\beamgain}{+5.90}
\newcommand{\greedyallgap}{−1.62}
\newcommand{\greedycopy}{65.3}
\newcommand{\humancopy}{48.8}
\newcommand{\greedyuntouched}{45.3}
\newcommand{\greedyoveredit}{13.6}
\newcommand{\beamoveredit}{12.7}
\newcommand{\greedysingle}{75.4}
\newcommand{\humansingle}{49.5}
\newcommand{\greedyratio}{1.00}
\newcommand{\strippedkept}{94.9}
\newcommand{\rawkept}{78.2}
\newcommand{\strippedlost}{1.4}
\newcommand{\rawlost}{17.5}
\newcommand{\strippedbert}{0.98}
\newcommand{\rawbert}{0.91}
\newcommand{\strippeddrop}{−0.04}
\newcommand{\rawdrop}{−0.45}
\newcommand{\strippedcopy}{68.6\%}
\newcommand{\rawcopy}{37.5\%}
\newcommand{\daasisari}{43.87}
\newcommand{\daasicopysari}{45.05}
\newcommand{\daasigain}{−1.19}
\newcommand{\daasibert}{0.99}
\newcommand{\daasidrop}{−0.03}
\newcommand{\daasicopy}{77.7\%}
\newcommand{\daasikept}{97.1}
\newcommand{\diacn}{3,918}
\newcommand{\diacrawkept}{63.7}
\newcommand{\diacrawlost}{31.9}
\newcommand{\diacstrippedkept}{94.6}
\newcommand{\diacstrippedlost}{1.8}
\newcommand{\resultstable}{%
\begin{table}[h]
\small
\caption{Model~1 on test data with human references, scored once with
\texttt{score.py}. SARI is split by whether the human reference changed the
source (DAASI has no unchanged pairs). BERTScore compares output with source;
readability drop is the fall in predicted BAREC level (positive = easier).
DAASI sources and references have tashkeel stripped.}
\label{tab:test}
\begin{tabularx}{\linewidth}{|L|C{11mm}|C{13mm}|C{15mm}|C{10mm}|C{13mm}|C{15mm}|C{11mm}|}
\hline
\head{System} & \head{SARI all} & \head{SARI changed} & \head{SARI un\-changed} & \head{BLEU} & \head{BERT\-Score} & \head{Read\-ability drop} & \head{Copy rate} \\ \hline
\multicolumn{8}{|l|}{\cellcolor{sichead}\textbf{SAMER test (3,277 rows, locked, in domain)}} \\ \hline
Copy the input & 77.50 & 56.06 & 100.00 & 72.03 & 1.00 & 0.00 & 100.0\% \\ \hline
Model~1, greedy & 75.88 & \textbf{61.83} & 90.62 & 71.68 & 0.98 & 0.08 & 65.3\% \\ \hline
Model~1, beam 4 & 76.29 & 61.95 & 91.34 & 72.22 & 0.98 & 0.07 & 69.4\% \\ \hline
\multicolumn{8}{|l|}{\cellcolor{sichead}\textbf{DAASI held-out (350 pairs, government and insurance text, out of domain)}} \\ \hline
Copy the input & 45.05 & 45.05 & --- & 25.25 & 1.00 & 0.00 & 100.0\% \\ \hline
Model~1, greedy & 43.87 & 43.87 & --- & 23.66 & 0.99 & −0.03 & 77.7\% \\ \hline
\end{tabularx}
\end{table}}
\newcommand{\barectable}{%
\begin{table}[h]
\small
\caption{Model~1 on the locked BAREC test set (7,286 sentences, no reference
simplifications), greedy decoding. Stripped: tashkeel removed from the source
before the model sees it, as the application does.}
\label{tab:barec}
\begin{tabularx}{\linewidth}{|L|C{17mm}|C{17mm}|C{17mm}|C{17mm}|C{17mm}|}
\hline
\head{Input} & \head{Copy rate} & \head{BERTScore} & \head{Read\-ability drop} & \head{Source words kept} & \head{Rows losing $>$half} \\ \hline
Stripped (as deployed) & 68.6\% & 0.98 & −0.04 & 94.9\% & 1.4\% \\ \hline
Raw, with tashkeel & 37.5\% & 0.91 & −0.45 & 78.2\% & 17.5\% \\ \hline
\end{tabularx}
\end{table}}
\newcommand{\figsari}{%
\begin{tikzpicture}
\begin{axis}[
  width=\linewidth, height=54mm, font=\footnotesize,
  ybar, bar width=13pt, ymin=40, ymax=112, ylabel={SARI},
  symbolic x coords={all rows, changed rows, unchanged rows},
  xtick=data, enlarge x limits=0.25,
  nodes near coords, nodes near coords style={font=\tiny, rotate=90, anchor=west,
    /pgf/number format/fixed, /pgf/number format/precision=1},
  legend style={at={(0.5,-0.2)}, anchor=north, legend columns=3,
                draw=black!25, font=\scriptsize},
  grid=major, grid style={black!8}, axis x line*=bottom, axis y line*=left]
  \addplot[draw=figgrey, fill=figgrey!45] coordinates {(all rows,77.50) (changed rows,56.06) (unchanged rows,100.00)};
  \addlegendentry{copy the input}
  \addplot[draw=sicblue, fill=sicblue!75] coordinates {(all rows,75.88) (changed rows,61.83) (unchanged rows,90.62)};
  \addlegendentry{model 1, greedy}
  \addplot[draw=figwarm, fill=figwarm!70] coordinates {(all rows,76.29) (changed rows,61.95) (unchanged rows,91.34)};
  \addlegendentry{model 1, beam 4}
\end{axis}
\end{tikzpicture}}
\newcommand{\figbehaviour}{%
\begin{tikzpicture}
\begin{axis}[
  name=a, width=0.42\linewidth, height=50mm, font=\scriptsize,
  xbar, bar width=6pt, xmin=0, xmax=100, xlabel={\%},
  symbolic y coords={single-word edits, changed rows left untouched, rows returned unchanged},
  ytick=data, enlarge y limits=0.3, yticklabel style={align=right, text width=19mm},
  nodes near coords, nodes near coords style={font=\tiny, /pgf/number format/fixed,
    /pgf/number format/precision=1},
  legend style={at={(0.5,-0.3)}, anchor=north, legend columns=2, draw=black!25},
  title={(a) SAMER test}, title style={font=\scriptsize\bfseries},
  grid=major, grid style={black!8}, axis x line*=bottom, axis y line*=left]
  \addplot[draw=sicblue, fill=sicblue!75] coordinates {(65.3,rows returned unchanged) (45.3,changed rows left untouched) (75.4,single-word edits)};
  \addlegendentry{model 1}
  \addplot[draw=figgrey, fill=figgrey!45] coordinates {(48.8,rows returned unchanged) (0,changed rows left untouched) (49.5,single-word edits)};
  \addlegendentry{human}
\end{axis}
\begin{axis}[
  at={(a.east)}, anchor=west, xshift=20mm,
  width=0.42\linewidth, height=50mm, font=\scriptsize,
  xbar, bar width=6pt, xmin=0, xmax=100, xlabel={\%},
  symbolic y coords={rows losing $>$half, source words kept},
  ytick=data, enlarge y limits=0.55, yticklabel style={align=right, text width=16mm},
  nodes near coords, nodes near coords style={font=\tiny, /pgf/number format/fixed,
    /pgf/number format/precision=1},
  legend style={at={(0.5,-0.3)}, anchor=north, legend columns=2, draw=black!25},
  title={(b) BAREC test, sentences with tashkeel}, title style={font=\scriptsize\bfseries},
  grid=major, grid style={black!8}, axis x line*=bottom, axis y line*=left]
  \addplot[draw=sicblue, fill=sicblue!75] coordinates {(94.6,source words kept) (1.8,rows losing $>$half)};
  \addlegendentry{tashkeel stripped}
  \addplot[draw=figwarm, fill=figwarm!75] coordinates {(63.7,source words kept) (31.9,rows losing $>$half)};
  \addlegendentry{raw input}
\end{axis}
\end{tikzpicture}}


\begin{document}

\siccover
  {Bayan: On-Device Arabic Text Simplification\\for Readers with Dyslexia\\[1mm]
   {\normalfont\fontsize{12}{15}\selectfont \ar{تبسيط}\enspace$\cdot$\enspace Simplify}}
  {23/09/26}
  {Cogni}
  {Marwan Elamami\\ Abdulrahman Khengari\\ Ahmed Alaeb\\ Sanad Ali\\
   Mohammed Thabet\\ Abdul Majid Mraied}

\siccontents

% ============================================================== 1 Introduction
\section{Introduction}

\subsection{Background Information}

Dyslexia impairs reading through phonological processing, and Modern Standard
Arabic adds three obstacles that tools built for Latin scripts do not address.
Short vowels are normally left unwritten, so one written form carries several
readings: \ar{كتب} is \emph{kataba} (he wrote), \emph{kutiba} (it was written)
or \emph{kutub} (books), and the reader must resolve it from context. Clitics
attach directly to the stem, producing long words that resist chunking. Prose
chains clauses with \ar{و} and \ar{ف} across whole paragraphs, so one sentence
can exceed working memory.

Adding every diacritic is not a complete answer; it adds visual clutter of its
own. The intervention we pursue is to \emph{simplify first, then vocalise}.
When the project started we found no deployed Arabic dyslexia-support tool and
no large, openly reusable Arabic simplification corpus. That shaped both what
we built and how much of our effort went into data and evaluation.

\subsection{Motivation and Objective}
\label{sec:objective}

Readers meet difficult Arabic inside other applications --- a message, a news
page, a PDF --- not in a reading app. Bayan therefore appears where the text
already is: selecting text on an Android phone offers \ar{تبسيط}
(\emph{Simplify}) in the system text-selection menu, and the simplified text
opens in a sheet over the same app. Every model runs on the phone, so the
reader's text never leaves it and the tool works offline.

We committed to four measurable objectives. \autoref{tab:objectives} in
\autoref{sec:accomplish} reports where each one stands today.

\begin{enumerate}[label=\textbf{O\arabic*}, leftmargin=8mm]
\item A simplifier that beats \emph{copying the input} on the sentences a human
      actually simplified, on held-out test data it never saw.
\item Diacritics restored on the simplified text without altering its letters.
\item The whole pipeline running on a mid-range Android phone, within 250\,MB for
      the simplifier.
\item Evidence from readers: a human evaluation with Arabic readers who have
      dyslexia, with informed consent.
\end{enumerate}

\subsection{Members and Role Assignments}

Every member has merged work into the team repository; the evidence column cites
the pull requests (\#) and issues that hold it.

\begin{table}[h]
\small
\caption{Members, roles, and the work each has delivered so far.}
\label{tab:roles}
\begin{tabularx}{\linewidth}{|p{27mm}|p{30mm}|L|}
\hline
\head{Member} & \head{Role} & \head{Delivered (evidence)} \\ \hline
Marwan Elamami & Team lead; integration; on-device model &
  Product-form decision (\#9); evaluation protocol and fixes (\#16, \#18);
  model~1 diagnosis; Baseet acquisition and leakage filtering; Action Plan;
  this report \\ \hline
Abdulrahman Khengari & Data lead &
  Validators for generated pairs (\#10, \#11); CAMeL readability check (\#17,
  \#20); readability classifier; synthetic corpus in progress (\#8, \#12) \\ \hline
Ahmed Alaeb & Model training &
  SAMER conversion; model~1 trained and documented (\#5, \#22) \\ \hline
Sanad Ali & Diacritization; co-training &
  Four-system diacritization benchmark; Libtashkeel selected (\#6, \#23) \\ \hline
Mohammed Thabet & Android application &
  \texttt{PROCESS\_TEXT} plugin, result sheet, 41 tests (\#14, \#21) \\ \hline
Abdul Majid Mraied & Evaluation &
  Scoring and leakage scripts (\#7, \#15); locked-set converter, copy
  baselines, generate-and-score path (\#19, \#24, \#26) \\ \hline
\end{tabularx}
\end{table}

\subsection{Schedule and Milestones}

The project runs from 13 September to the pitch on 6 October 2026. At the time
of writing (23 September, day 10 of 23) the first build-and-measure cycle is
complete: a model has been trained, scored against a baseline on locked test
sets, and diagnosed, and the application shell exists. \autoref{fig:schedule}
shows what was done, what is running, and what remains; \autoref{tab:milestones}
compares the dates we set with the dates we met.

\begin{figure}[h]
\centering
\figschedule
\caption{Project timeline as executed. Blue: done. Orange: in progress, solid to
today. Dashed: not started. AP = Action Plan; M1, M2 = mentoring sessions.
Numbers after a task are its GitHub issues.}
\label{fig:schedule}
\end{figure}

\begin{table}[h]
\small
\caption{Milestones: target against actual.}
\label{tab:milestones}
\begin{tabularx}{\linewidth}{|L|C{19mm}|C{19mm}|C{21mm}|}
\hline
\head{Milestone} & \head{Target} & \head{Actual} & \head{Status} \\ \hline
Evaluation scripts ready (\#7)                 & Mon 21 Sep & Mon 21 Sep & \done \\ \hline
Diacritizer selected (\#6)                     & Mon 21 Sep & Mon 21 Sep & \done \\ \hline
Android plugin shell (\#14)                    & ---        & Mon 21 Sep & \done \\ \hline
Action Plan submitted                          & Tue 22 Sep & Tue 22 Sep & \done \\ \hline
Model~1 scored on the locked test sets         & Wed 23 Sep & Wed 23 Sep & \done \\ \hline
Synthetic corpus exported (\#8)                & Fri 25 Sep & ---        & \wip \\ \hline
Model~2 trained and scored                     & Tue 29 Sep & ---        & \wip \\ \hline
Model running on a phone, size and latency measured & Fri 2 Oct & ---   & \todo \\ \hline
Human evaluation with readers                  & Sat 3 Oct  & ---        & \todo \\ \hline
Final submission / pitch                       & 5 / 6 Oct  & ---        & \todo \\ \hline
\end{tabularx}
\end{table}

% ========================================================= 2 Project Execution
\section{Project Execution}

\subsection{Data Acquisition}

We acquired five corpora (\autoref{tab:corpora}). SAMER is used under written
approval from CAMeL Lab to fine-tune and to publish the resulting weights
non-commercially; the corpus itself is never redistributed, and no SAMER sentence
is reproduced in this report. Baseet is used with its authors' permission. Two test
splits --- SAMER test and BAREC test --- were locked before any training and
are never trained on.

\begin{table}[h]
\small
\caption{Corpora acquired. Sizes are rows in the files we hold.}
\label{tab:corpora}
\begin{tabularx}{\linewidth}{|p{22mm}|p{33mm}|p{22mm}|L|}
\hline
\head{Corpus} & \head{Size} & \head{Licence} & \head{Used for} \\ \hline
SAMER \cite{samer} & 14,343 / 2,983 / 3,277 pairs (train/dev/test) &
  research; written approval & model~1 training; locked test set \\ \hline
BAREC \cite{barec} & 54,845 / 7,310 / 7,286 sentences & CC BY-SA 4.0 &
  out-of-domain locked test set; readability labels \\ \hline
Baseet \cite{baseet} & 34,294 sources, three levels each; 31,104 after filtering &
  with permission & general-domain training pairs for model~2 \\ \hline
DAASI \cite{daasi} & 1,509 pairs + 350 held out & CC0 &
  human-written pairs; out-of-domain check \\ \hline
Tashkeela \cite{tashkeela} & 743 test sentences & open &
  diacritization benchmark \\ \hline
\end{tabularx}
\end{table}

A sixth source, simplification pairs generated by a large language model and
admitted only through validators and human review, is being produced
(\#8, \#12). The validators are merged (\#11, \#20); the corpus has not yet been
delivered, and nothing in this report was trained on it.

\subsection{Training Methodology}
\label{sec:method}

\lead{Candidate base models.} Simplification is a sequence-to-sequence task,
so every candidate is a pretrained encoder--decoder transformer: an encoder reads
the whole source sentence, and a decoder writes the simplification one token at
a time while attending to it. We short-listed three Arabic models of this kind
(\autoref{tab:basemodels}). Architecture figures are read from each model's
configuration file; parameter counts from its fp32 checkpoint size.

\begin{table}[h]
\small
\caption{Candidate base models. ``Tashkeel'' is measured on 1,000 diacritized
BAREC dev sentences: vocabulary entries containing a diacritic, and how many
more tokens the same sentences cost with diacritics than without.}
\label{tab:basemodels}
\begin{tabularx}{\linewidth}{|p{27mm}|L|L|L|}
\hline
\head{} & \head{AraT5v2-base-1024} & \head{AraBART} & \head{HPLT T5 base (Arabic)} \\ \hline
Family & T5 \cite{t5}: span-corruption pretraining, relative position buckets &
  BART \cite{bart}: denoising pretraining, learned positions &
  T5 in the NorT5 variant \cite{nort5} \\ \hline
Layers (enc/dec), width & 12/12, $d$=768, 12 heads, gated-GELU FFN 2048 &
  6/6, $d$=768, 12 heads, GELU FFN 3072 &
  12/12, $d$=768, 12 heads, FFN 2048 \\ \hline
Parameters & 368M (46\% in two untied 110k$\times$768 embeddings) & 139M & $\approx$294M \\ \hline
Vocabulary & 110,208 SentencePiece & 50,002 SentencePiece BPE & 32,768 \\ \hline
Context & 1,024 tokens & 1,024 positions & 512 positions \\ \hline
Pretraining & Arabic MSA and dialects; v2 of AraT5 \cite{arat5,octopus} &
  large Arabic corpora \cite{arabart} & HPLT v3 web crawl, Arabic subset \cite{hplt} \\ \hline
Tashkeel & 1.6\% of pieces; 2.57$\times$ tokens &
  0 pieces; every mark becomes \texttt{<unk>}; 4.53$\times$ tokens &
  0 pieces; 1.81$\times$ tokens \\ \hline
Licence & none stated on the model card & Apache-2.0 & Apache-2.0 \\ \hline
Status & model~1 trained & queued for model~2 & held in reserve \\ \hline
\end{tabularx}
\end{table}

AraT5v2 went first because, on Baseet's head-to-head evaluation (same split,
same harness), it beat AraBART at all three simplification levels. AraBART is
the challenger because it is 2.6 times smaller, which makes decoding on a phone
cheaper. HPLT T5 is the reserve: it is the most recent, but needs custom loading
code and has the shortest context.

\lead{Why the input is stripped of tashkeel.} None of the three vocabularies was
built for diacritized text. In AraT5v2 only 1.6\% of pieces carry a diacritic,
and a diacritized sentence costs 2.57 times as many tokens; AraBART cannot
represent a diacritic at all and turns each one into an unknown token. The
models learned Arabic from mostly undiacritized text and resolve ambiguous words
from context, so a diacritic adds little they can use and a great deal they
cannot. For every model trained on SAMER, which has no diacritics, the source is
stripped before it reaches the simplifier; the diacritizer restores vowels
afterwards, on the simplified text.

\lead{Model~1.} We fine-tuned AraT5v2-base-1024 \cite{arat5} (368M parameters),
a text-to-text transformer pretrained on Arabic, to rewrite SAMER's level-5
sentences into their level-3 simplifications. Training used AdamW (learning
rate $10^{-4}$, weight decay 0.01, 100 warm-up steps), an effective batch of 16,
mixed precision and gradient checkpointing, for 10 epochs (3,250 steps) with
seed 42. Pairs whose simplification is identical to the source make up 42\% of
SAMER train; they were capped at 20\% of training so the model sees mostly real
rewrites, and left untouched in dev and test. Each source was given the task prefix
\ar{بسّط:}\ (\emph{simplify:}). Every 200 steps SARI was computed
on a fixed 500-row dev subset, and the checkpoint with the best SARI (step
3,000) was kept.

\lead{How it is evaluated.} SARI \cite{sari} is the main metric, computed with
the \texttt{evaluate} implementation, alongside BLEU, BERTScore \cite{bertscore}
between source and output (meaning preservation), and the drop in predicted
readability level. Every score is reported beside the \emph{copy baseline} ---
the same metric for an output identical to the input --- and split into rows a
human changed and rows a human left unchanged (\autoref{sec:modeling} shows why
the split matters). Model~1 was scored once on each locked test set after
training was finished; no training or selection decision used those scores.

\subsection{Workflow}

Work runs through GitHub: each task is an issue with an owner, a reviewer and a
definition of done; changes arrive as pull requests, are reviewed, and are
squash-merged. By 23 September the team had opened 12 issues (9 closed) and
merged 12 pull requests, with at least one from every member.

\autoref{fig:pipeline} shows the data path. Two gates can reject data. The
leakage check compares every training file against both locked test sets, by
exact match and by character 5-gram Jaccard overlap of 0.8 or more; its first
real use removed 9.3\% of a published corpus (\autoref{sec:prep}). Generated
pairs pass validators and human review before they may be trained on. The locked
test sets are fixed by row count and SHA-256 in a committed manifest, and both
hashes were re-verified before the test runs reported here.

\begin{figure}[h]
\centering
\figpipeline
\caption{Data and evaluation workflow. Orange boxes are the gates that can
reject data. Solid boxes are running; the dashed branch on the right (generated
data) is built but has not yet delivered a corpus.}
\label{fig:pipeline}
\end{figure}

\subsection{System Design}

\autoref{fig:arch} shows the system and the state of each part. The Android
application registers an activity for \texttt{android.intent.action.PROCESS\_TEXT}
(mime type \texttt{text/plain}, Android 6.0 and later). A translucent theme lets
the result sheet float over the host app, so the \texttt{SYSTEM\_ALERT\_WINDOW}
overlay permission is never requested. Inference sits behind a
\texttt{TextSimplifier} interface; today a mock engine implements it, so the
interface, sheet and tests are complete while the model is integrated.

On the phone, models run through ONNX Runtime Mobile with greedy decoding and
the tokenizer exported inside the model graph, so no Python runtime is needed.
Diacritics are stripped before the simplifier and restored after it, for the
reason measured in \autoref{sec:modeling}. The simplifier's size budget is
250\,MB. AraT5v2 is 46\% embedding matrices, so vocabulary pruning plus int8
quantization is estimated at about 223\,MB; that is a calculation, and the
on-device measurement is scheduled (\autoref{sec:future}). The selected
diacritizer is 4.8\,MB, measured.

\begin{figure}[h]
\centering
\figarchitecture
\caption{System architecture and build status on 23 September. The plugin is
entered from the host app's selection menu and answers with a sheet over it.}
\label{fig:arch}
\end{figure}

% ================================================================== 3 Results
\section{Results}

\subsection{Data Preprocessing}
\label{sec:prep}

\lead{Conversion.} SAMER's aligned levels were converted to source--target pairs
(level~5 $\rightarrow$ level~3). A single converter
(\texttt{make\_eval\_input.py}) turns each locked test file into the scoring
format, keeps every row --- including the 1,599 SAMER test rows (48.8\%) where
the two levels are identical --- and checks the row count before writing.

\lead{Normalisation for comparison only.} For leakage checks and scoring
alignment, diacritics and punctuation are removed and alef, \ar{ة}/\ar{ه} and
\ar{ى}/\ar{ي} variants are unified. Training text keeps its original spelling.

\lead{Leakage checks.} The official splits, model~1's training file and a sample
of generated data were all checked against both locked test sets and found
clean. Matches under five words are reported but not counted, because short
formulaic fragments recur in any corpus; BAREC alone repeats 311 sentences
between its own train and test splits.

\finding{\textbf{Contamination found and removed.} 3,190 of Baseet's 34,294 rows
(9.3\%) overlapped our locked test sets --- 1,829 SAMER test sentences and 1,361
BAREC test sentences. Seen the other way, 47\% of SAMER's substantive test
sentences appear somewhere in Baseet. Training on it unfiltered would have
inflated every later score without any visible error. The rows were removed and
the remaining 31,104 re-checked: zero matches.}

\subsection{Exploratory Data Analysis (EDA)}

We measured the corpora rather than assume them (\autoref{tab:eda},
\autoref{fig:lengths}); three differences turned out to explain model~1's
behaviour.

\begin{table}[h]
\small
\caption{Corpus statistics, computed on the files we hold. Identity = target
equals source after removing diacritics; length ratio = target words over source
words, median.}
\label{tab:eda}
\begin{tabularx}{\linewidth}{|L|C{17mm}|C{15mm}|C{15mm}|C{17mm}|C{15mm}|C{15mm}|}
\hline
\head{Corpus} & \head{Rows} & \head{Median words} & \head{$\le$3 words} & \head{Diacritized} & \head{Identity} & \head{Length ratio} \\ \hline
SAMER train          & 14,343 &  7 & 16.4\% & 0.0\%  & 42.1\% & 1.00 \\ \hline
SAMER test (locked)  &  3,277 &  6 & 23.6\% & 0.0\%  & 48.8\% & 1.00 \\ \hline
BAREC test (locked)  &  7,286 & 10 & 16.9\% & 53.8\% & ---    & ---  \\ \hline
Baseet, strongest level & 31,104 & 14 & 0.0\% & 36.5\% & 0.1\% & 0.73 \\ \hline
DAASI                &  1,509 & 12 &  8.4\% & 32.3\% & 0.0\%  & 1.00 \\ \hline
\end{tabularx}
\end{table}

\begin{figure}[h]
\centering
\figlengths
\caption{Source sentence length by corpus. SAMER is short literary text; Baseet
and BAREC reach the long sentences a reader actually struggles with.}
\label{fig:lengths}
\end{figure}

\begin{itemize}
\item \textbf{SAMER has no diacritics at all}, while 54\% of BAREC test and 36\%
      of Baseet carry them. A model trained only on SAMER has never seen a
      diacritized input.
\item \textbf{SAMER simplifies words, not sentences.} Its median length ratio is
      exactly 1.00: humans substituted vocabulary and kept structure. Baseet's
      strongest level shortens a 14-word median source to 9 words (ratio 0.73).
\item \textbf{Half of SAMER needs no change.} 42--49\% of pairs are identical,
      and SAMER splits by chapter rather than by novel, so its test set is fully
      in-domain. Both make SARI on SAMER read optimistically.
\end{itemize}

\subsection{Modeling}
\label{sec:modeling}

\lead{The baseline comes first.} Because half of SAMER needs no change, copying
the input scores well: SARI \copytestall{} on the full test set, but only
\copytestchanged{} on the \testchangedn{} rows a human changed. The second number
is the bar a model must clear. BAREC has no reference simplifications, so on it
only the reference-free measures apply.

\lead{Model~1 on the locked test sets.} We ran model~1 once on each locked
test set, with the task prefix it was trained on, under greedy decoding (what
the phone will run) and under beam search with four beams
(\autoref{tab:test}, \autoref{fig:sari}). Before the test run the harness was
validated on dev: it reproduced Ahmed's 2,983 dev predictions exactly, row for
row.

\resultstable

\begin{figure}[h]
\centering
\figsari
\caption{SARI on the locked SAMER test set by reference bucket. Model~1 clears
the copy baseline where a human simplified, and loses to it where a human left
the sentence alone.}
\label{fig:sari}
\end{figure}

\finding{\textbf{Model~1 simplifies where it should, and over-edits where it
should not.} On the \testchangedn{} test rows a human simplified, greedy decoding
scores \greedychanged{} against the copy baseline's \copytestchanged{}
(\greedygain{} SARI): the model is doing real work. On the
\testunchangedn{} rows a human left alone it scores \greedyunchanged{} against a
perfect 100, because it edits \greedyoveredit\% of sentences that needed no
change. The two effects net to \greedyallgap{} overall. Beam search does not
change the picture (\beamchanged{} on changed rows), so greedy decoding costs
nothing on the phone.}

\lead{It learned what SAMER teaches.} Against the human references on test, the
model returns \greedycopy\% of rows unchanged where humans left
\humancopy\%, and leaves \greedyuntouched\% of the rows a human rewrote
untouched. When it edits, \greedysingle\% of its edits change a single word
(humans: \humansingle\%), and its median output-to-input length ratio is
\greedyratio: it substitutes vocabulary and never restructures
(\autoref{fig:behaviour}a), exactly as the EDA predicts for SAMER. The
readability classifier agrees: outputs are on average 0.08 of a level easier
than their sources on BAREC's 19-level scale, a change a reader would barely
notice.
\autoref{fig:training} shows how checkpoint selection contributed: copy rate
was still falling at the last step, while the SARI that chose the checkpoint
had been flat within noise since step 1,800, and validation loss had bottomed
at step 2,200.

\begin{figure}[h]
\centering
\figtraining
\caption{Model~1 training. The checkpoint was selected on a 0.18-point SARI
difference while the model's copying behaviour was still changing.}
\label{fig:training}
\end{figure}

\lead{Out of domain, stripping tashkeel is what keeps it usable.} BAREC's
7,286 test sentences are general-domain text, and \diacn{} of them carry
tashkeel. We ran the model twice: on the raw text, and with tashkeel stripped
from the source as the application does (\autoref{tab:barec},
\autoref{fig:behaviour}b). On the sentences that carry tashkeel, raw input keeps
only \diacrawkept\% of the source words and \diacrawlost\% of outputs lose more
than half the sentence; stripped, the same sentences keep \diacstrippedkept\% and
only \diacstrippedlost\% lose more than half. Stripping removes the model's worst
failure at no measurable cost.

\lead{Out of domain, it does not yet beat copying.} DAASI's held-out split is
government and insurance text with human simplifications, so it gives an
out-of-domain SARI. There model~1 scores \daasisari{} against a copy baseline of
\daasicopysari{} (\daasigain{}): it returns \daasicopy{} of sentences unchanged,
having learned from SAMER's novels when \emph{not} to edit, not how to simplify
administrative prose. This is the gap model~2's general-domain training data
is chosen to close (\autoref{tab:improvements}).

\barectable

\begin{figure}[h]
\centering
\figbehaviour
\caption{Model~1 behaviour. (a) SAMER test, greedy, against the human reference.
(b) BAREC test, the \diacn{} sentences that carry tashkeel, greedy, with and
without tashkeel stripped from the input.}
\label{fig:behaviour}
\end{figure}

\lead{Meaning errors are not rare.} On dev, 832 of the model's substitutions
disagree with the human, 347 of them on words the human deliberately kept.
Confirmed examples: \ar{عزلة} (isolation) $\rightarrow$ \ar{استعباد}
(enslavement), \ar{ملجم} (bridled) $\rightarrow$ \ar{مسلح} (armed). SARI and
BLEU do not detect this class of error, which is why human error annotation is
part of the evaluation.

\lead{Diacritization.} Four systems were benchmarked on the 743-sentence
Tashkeela test fixture (\autoref{tab:diac}). Libtashkeel was selected: CATT is
statistically tied on accuracy (0.01 points, about six characters in 67,576) but
16 times larger and carries a non-commercial licence. CAMeL Tools was rejected
for a stronger reason than accuracy: it changed the base letters of 32 of 100
sentences, corrupting the text it was asked to vocalise.

\begin{table}[h]
\small
\caption{Diacritization benchmark (\#23). DER = diacritic error rate; base-letter
corruption measured on 100 sentences.}
\label{tab:diac}
\begin{tabularx}{\linewidth}{|L|C{13mm}|C{15mm}|C{13mm}|C{18mm}|C{14mm}|C{16mm}|}
\hline
\head{System} & \head{DER} & \head{DER, no case} & \head{WER} & \head{Base letters changed} & \head{Size} & \head{Licence} \\ \hline
CAMeL Tools MLE \cite{cameltools} & 16.57\% & 13.97\% & 56.73\% & \textbf{32\%} & 129\,MB & MIT/GPL \\ \hline
CATT \cite{catt}                  & 8.44\%  & 6.92\%  & 29.40\% & 0\% & 77.9\,MB & CC BY-NC \\ \hline
Mishkal \cite{mishkal}            & 12.74\% & 10.12\% & 39.33\% & 0\% & 4.3\,MB & GPLv3 \\ \hline
\textbf{Libtashkeel} \cite{libtashkeel} & \textbf{8.43\%} & \textbf{6.91\%} & 29.57\% & 0\% & \textbf{4.8\,MB} & MIT \\ \hline
\end{tabularx}
\end{table}

\subsection{User Interface}

\autoref{fig:ui} shows the interface as built. Selecting Arabic text in another
app offers \ar{تبسيط}; choosing it opens a Material bottom sheet over the app
with the simplified text first, the original below it in a muted card, and
controls for listening, copying and settings. Typography follows dyslexia
guidance rather than density: 20\,sp text, line height 1.6 plus 10\,sp, letter
spacing 0.04, and right-to-left text left \emph{unjustified}, because
justification in Arabic stretches words and opens uneven gaps. Swiping down,
tapping outside or pressing Back closes the sheet and returns to the host app.

We documented the expected selection behaviour in five host apps --- Chrome,
WhatsApp, Keep, PDF readers and Gmail --- including where \ar{تبسيط} appears in
each menu and which inputs are read-only. Confirming this on physical devices is
part of the on-device milestone. Listen and copy are visible but disabled until
the read-aloud and model engines are connected.

\begin{figure}[h]
\centering
\figui
\caption{The interface as implemented (\#21), drawn from its layout file. The
sheet's text is a placeholder: the engine behind it is still the mock.}
\label{fig:ui}
\end{figure}

\subsection{Testing and Improvements}

\lead{Testing.} The application ships 41 unit and Robolectric tests covering
intent extraction, empty input, view binding and the typography attributes. The
leakage checker is tested against a planted duplicate and must exit non-zero.
Running the scoring script end to end on real data, rather than fixtures,
exposed a bug that put every BAREC row in the ``unchanged'' bucket; BAREC now
has its own reference-free bucket. The same end-to-end run showed that the
untrained AraT5v2 copies almost nothing and scores SARI 22.6 on 200 dev rows,
so the gain from fine-tuning is real even though model~1 is conservative.
Preparing the test run exposed a second harness bug: model~1 was trained with
the input prefix \ar{بسّط:}, and the shared prediction script did not add it.
Without the prefix the model inserts spurious words and scores 65.2 SARI on the
test set, twelve points below copying; with it, the harness reproduces the
training run's predictions exactly.

\lead{Improvements.} Each measured problem produced a change
(\autoref{tab:improvements}).

\begin{table}[h]
\small
\caption{What the measurements changed.}
\label{tab:improvements}
\begin{tabularx}{\linewidth}{|L|L|C{21mm}|}
\hline
\head{Finding} & \head{Change} & \head{Status} \\ \hline
A copying model scores SARI 77.5 on SAMER test &
  Copy baseline and changed/unchanged split reported with every score & \done \\ \hline
9.3\% of Baseet overlaps our test sets & Rows removed; leakage check required for every training file & \done \\ \hline
Model collapses on diacritized input & Strip diacritics before the simplifier; diacritized sources in training & \done{} / \wip \\ \hline
Checkpoint chosen on SARI noise, over copy-heavy dev & Stop on validation loss; select on changed-row SARI & \wip \\ \hline
Prediction script omitted the training prefix & Prefix made a required, recorded setting of every model & \wip \\ \hline
Model edits \greedyoveredit\% of sentences that need no change & Identity rows kept in training at a controlled share; leave-unchanged mode as a level token & \wip \\ \hline
SAMER teaches word swaps only; below copying on DAASI (\daasigain{}) & Model~2 corpus of $\approx$106,700 pairs: Baseet (three levels), SAMER, DAASI, identity rows capped at 4.6\% & \wip \\ \hline
Published generation config stops at 20 tokens; licence tag says Apache-2.0 & Fix \texttt{max\_length}; set the non-commercial licence SAMER requires & \todo \\ \hline
One architecture tried & AraBART \cite{arabart} (139M) trained on the same data as AraT5v2 & \todo \\ \hline
\end{tabularx}
\end{table}

% ========================================================== 4 Projected Impact
\section{Projected Impact}

\subsection{Accomplishments and Benefits}
\label{sec:accomplish}

\begin{table}[h]
\small
\caption{Objectives set in \autoref{sec:objective}, against the evidence on 23
September.}
\label{tab:objectives}
\begin{tabularx}{\linewidth}{|p{47mm}|L|C{21mm}|}
\hline
\head{Objective} & \head{Evidence so far} & \head{Status} \\ \hline
O1 Beat copying on changed rows, unseen data & Met in domain: SAMER test \greedychanged{} vs \copytestchanged{} (\greedygain{}). Not yet out of domain: DAASI \daasisari{} vs \daasicopysari{} (\daasigain{}) & \wip \\ \hline
O2 Diacritize without changing letters & Libtashkeel: DER 6.91\% without case endings, 0 of 100 sentences altered, 4.8\,MB & \done \\ \hline
O3 Run on a phone within 250\,MB & Plugin and sheet built and tested; model integration and measurement pending & \wip \\ \hline
O4 Evidence from readers & Protocol drafted; sessions not yet run & \todo \\ \hline
\end{tabularx}
\end{table}

\lead{Built so far.} A working Android plugin on the \texttt{PROCESS\_TEXT}
intent; a trained simplifier, scored on two locked test sets beside a copy
baseline; a selected diacritizer with a four-system benchmark behind it; a
hashed test-set manifest, a reusable leakage checker, and a scoring pipeline
any member can run from one command.

\lead{Benefits.} For readers: help at the moment and place the difficulty
occurs, offline and private, rather than in a separate app. For schools and
teachers: a tool that works on the phones students already own. For the field:
an evaluation protocol that exposes copy-heavy scores, and a documented case of
test-set contamination in a public Arabic simplification corpus.

\subsection{Future Improvements}
\label{sec:future}

\lead{Before submission (5 October).} Train model~2 on the mixed corpus with the
corrected selection rule, alongside AraBART, and choose on changed-row SARI;
export the winner to ONNX, quantize it to int8 and measure size and latency on a
real phone; connect diacritization and read-aloud; run the human evaluation with
Arabic readers who have dyslexia; report annotator agreement on the generated
data.

\lead{Known limitations.} Automatic metrics miss meaning-altering
substitutions, so human annotation stays necessary. SAMER's test set is
in-domain, so SARI on it overstates performance on news and forms; DAASI's
held-out split is our out-of-domain check.

\lead{Beyond the capstone.} Sentence splitting, which SAMER does not teach and
which matters most to readers with dyslexia; a diacritizer fine-tuned on
simplified text; more domains in the corpus.

\subsection*{References}
\addcontentsline{toc}{subsection}{References}
\begin{thebibliography}{99}
\small
\setlength{\itemsep}{1pt}
\bibitem{samer} B.~Alhafni, R.~Hazim, J.~D. Piñeros Liberato, M.~Al Khalil and N.~Habash. The SAMER Arabic Text Simplification Corpus. \emph{LREC-COLING}, pp.~16079--16093, 2024.
\bibitem{barec} K.~Elmadani, N.~Habash and H.~Taha. A Large and Balanced Corpus for Fine-grained Arabic Readability Assessment. \emph{Findings of ACL}, pp.~16376--16400, 2025.
\bibitem{baseet} Baseet: a multi-level dataset for Arabic text simplification. \url{https://github.com/raya-y/Baseet}
\bibitem{daasi} DAASI: Arabic simplification pairs for government and insurance text. CC0.
\bibitem{tashkeela} T.~Zerrouki and A.~Balla. Tashkeela: Novel corpus of Arabic vocalized texts, data for auto-diacritization systems. \emph{Data in Brief}, 11:147--151, 2017.
\bibitem{arat5} E.~M.~B. Nagoudi, A.~Elmadany and M.~Abdul-Mageed. AraT5: Text-to-Text Transformers for Arabic Language Generation. \emph{ACL}, pp.~628--647, 2022.
\bibitem{arabart} M.~K. Eddine, N.~Tomeh, N.~Habash, J.~Le Roux and M.~Vazirgiannis. AraBART: a Pretrained Arabic Sequence-to-Sequence Model for Abstractive Summarization. \emph{WANLP}, pp.~31--42, 2022.
\bibitem{octopus} A.~Elmadany, E.~M.~B. Nagoudi and M.~Abdul-Mageed. Octopus: A Multitask Model and Toolkit for Arabic Natural Language Generation. \emph{ArabicNLP}, pp.~232--243, 2023. (Introduces AraT5v2.)
\bibitem{t5} C.~Raffel, N.~Shazeer, A.~Roberts, K.~Lee, S.~Narang, M.~Matena, Y.~Zhou, W.~Li and P.~J. Liu. Exploring the Limits of Transfer Learning with a Unified Text-to-Text Transformer. \emph{JMLR}, 21(140):1--67, 2020.
\bibitem{bart} M.~Lewis, Y.~Liu, N.~Goyal, M.~Ghazvininejad, A.~Mohamed, O.~Levy, V.~Stoyanov and L.~Zettlemoyer. BART: Denoising Sequence-to-Sequence Pre-training for Natural Language Generation, Translation, and Comprehension. \emph{ACL}, pp.~7871--7880, 2020.
\bibitem{nort5} D.~Samuel, A.~Kutuzov, S.~Touileb, E.~Velldal, L.~Øvrelid, E.~Rønningstad, E.~Sigdel and A.~Palatkina. NorBench -- A Benchmark for Norwegian Language Models. \emph{NoDaLiDa}, pp.~618--633, 2023.
\bibitem{hplt} L.~Burchell et al. An Expanded Massive Multilingual Dataset for High-Performance Language Technologies (HPLT). \emph{ACL}, pp.~17452--17485, 2025.
\bibitem{sari} W.~Xu, C.~Napoles, E.~Pavlick, Q.~Chen and C.~Callison-Burch. Optimizing Statistical Machine Translation for Text Simplification. \emph{TACL}, 4:401--415, 2016.
\bibitem{bertscore} T.~Zhang, V.~Kishore, F.~Wu, K.~Q. Weinberger and Y.~Artzi. BERTScore: Evaluating Text Generation with BERT. \emph{ICLR}, 2020.
\bibitem{cameltools} O.~Obeid et al. CAMeL Tools: An Open Source Python Toolkit for Arabic Natural Language Processing. \emph{LREC}, pp.~7022--7032, 2020.
\bibitem{catt} abjadai. CATT: Character-based Arabic Tashkeel Transformer. \url{https://github.com/abjadai/catt}
\bibitem{mishkal} T.~Zerrouki. Mishkal: Arabic text vocalization. \url{https://github.com/linuxscout/mishkal}
\bibitem{libtashkeel} Libtashkeel, an ONNX Arabic diacritization model (MIT), via the \texttt{text2tashkeel} wrapper.
\end{thebibliography}

% ====================================================== 5 Team Member Review
\clearpage
\section{Team Member Review and Comment}

\begin{fullwidth}
\setlength{\fboxrule}{0.5pt}\fcolorbox{sicrule}{white}{\parbox[c][30mm][c]{\dimexpr\linewidth-2\fboxsep-1pt\relax}{%
  \centering\fontsize{14}{16}\selectfont\color{sicrule}<ATTACH A TEAM PICTURE HERE>}}

\vspace{3mm}
\renewcommand{\arraystretch}{1.9}
\begin{tabularx}{\linewidth}{|p{38mm}|X|}
\hline
\head{NAME} & \head{REVIEW and COMMENT} \\ \hline
Marwan Elamami & \\ \hline
Abdulrahman Khengari & \\ \hline
Ahmed Alaeb & \\ \hline
Sanad Ali & \\ \hline
Mohammed Thabet & \\ \hline
Abdul Majid Mraied & \\ \hline
\end{tabularx}
\end{fullwidth}

% =================================================== 6 Instructor Review
\section{Instructor Review and Comment}

\begin{fullwidth}
\renewcommand{\arraystretch}{1}
\newcommand{\cat}[1]{\parbox[c][23mm][c]{27mm}{\centering\bfseries\color{sicblue}#1}}
\newcommand{\score}[1]{\parbox[c][23mm][c]{26mm}{\centering #1}}
\begin{tabularx}{\linewidth}{|p{27mm}|p{26mm}|X|}
\hline
\head{CATEGORY} & \head{SCORE} & \head{REVIEW and COMMENT} \\ \hline
\cat{IDEA} & \score{\_\_/10} & \\ \hline
\cat{APPLICATION} & \score{\_\_/30} & \\ \hline
\cat{RESULT} & \score{\_\_/30} & \\ \hline
\cat{PROJECT\\MANAGEMENT} & \score{\_\_/10} & \\ \hline
\cat{PRESENTATION\\\& REPORT} & \score{\_\_/20} & \\ \hline
\cat{TOTAL} & \score{\_\_/100} & \\ \hline
\end{tabularx}
\end{fullwidth}

\end{document}
